%%%%%%%%%%%%%%%%%%%%%%%%%%%%%%%%%%%%%%%%%%%%%%%%%%%%%%%%%%%%%%%%%%%%%%
% Lớp 11 | Hàm số lượng giác và phương trình lượng giác | [Giảm] Phương trình lượng giác thường gặp
% Dạng 1D1?6-4: Phương trình đối xứng, phản đối xứng
%%%%%%%%%%%%%%%%%%%%%%%%%%%%%%%%%%%%%%%%%%%%%%%%%%%%%%%%%%%%%%%%%%%%%%


% nguồn: 2026-11-15_11_1D1_Luong_giac (nhập 14/11/2026)
\begin{bt}%[1D1V6-4]%[ma:3304086d]
% nguon: Archive/SP_KhoLatex/ID6_old/12-OT-PTLuongGiacThuongGap.tex #3700

	Giải phương trình $ \sin 2x+(2-\sqrt{2})(\sin x+\cos x)+1-2\sqrt{2}=0 $.
	\loigiai{
	%	$ \sin 2x+(2-\sqrt{2})(\sin x+\cos x)+1-2\sqrt{2}=0 $. \tagEX{1}
	\noindent	Đặt $ t=\sin x+\cos x $	($ |t|\leq \sqrt{2} $), suy ra $ t^{2}=1+2\sin x\cos x\Rightarrow \sin 2x=t^{2}-1 $.\\
	Phương trình đã cho trở thành
		% Thay vào phương trình (1) ta được
		\[
		t^{2}-1+(2-\sqrt{2})t+1-2\sqrt{2}=0\Leftrightarrow t^{2}+(2-\sqrt{2})t-2\sqrt{2}=0\Leftrightarrow\hoac{& t=\sqrt{2}\text{ (thỏa mãn)}\\ &t=-2\text{ (loại)}.}
		\]
		Với $ t=\sqrt{2} $, ta có
		$$ \sin x+\cos x=\sqrt{2}\Leftrightarrow \dfrac{1}{\sqrt{2}}\sin x+\dfrac{1}{\sqrt{2}}\cos x=1\Leftrightarrow \cos \left(x-\dfrac{\pi}{4}\right)=1\Leftrightarrow x=\dfrac{\pi}{4}+k2\pi, \left(k\in\mathbb{Z}\right).$$
	}
\end{bt}

% nguồn: 2026-11-15_11_1D1_Luong_giac (nhập 14/11/2026)
\begin{bt}%[1D1H6-4]%[ma:249f1278]
% nguon: Archive/SP_KhoLatex/ID6_old/12-OT-PTLuongGiacThuongGap.tex #3701

	Giải phương trình $ 2(\tan^{2}x+\cot^{2}x)-(4-\sqrt{2})(\tan x+\cot x)+4+2\sqrt{2}=0 $.
	\loigiai
	{
		Điều kiện xác định $ \heva{& \sin x\neq 0\\ &\cos x\neq 0}\Leftrightarrow x\neq k\dfrac{\pi}{2} $.\\
		Đặt $ t=\tan x+\cot x $	($ |t|\geq 2 $), suy ra $ t^{2}=\tan^{2}x+\cot^{2}x+2\Rightarrow \tan^{2}x+\cot^{2}x=t^{2}-2 $. Thay vào (1) ta được
		\[
		2(t^{2}-2)-(4-\sqrt{2})t+4+2\sqrt{2}=0\Leftrightarrow 2t^{2}-(4-\sqrt{2})t+2\sqrt{2}=0\text{ (vô nghiệm)}.
		\]
		Vậy phương trình vô nghiệm.
	}
\end{bt}

% nguồn: 2026-11-15_11_1D1_TL2_Luong_giac (nhập 15/11/2026)
\begin{bt}%[1D1V6-4]%[ma:4a537119]
% nguon: Archive/SP_KhoLatex/ID6_old/1K1-4-10.tex #3614
%[DCHT-11-KNTT]%[Vĩ Lê Văn]
	Giải phương trình $2(\sin x+\cos x)+\sin 2 x+1=0$.
	\dapso{$x=-\dfrac{\pi}{4}+k \pi, \quad k \in \mathbb{Z}$.}
	\loigiai{
		Ta có $2(\sin x+\cos x)+\sin 2 x+1=0
		\Leftrightarrow 2(\sin x+\cos x)+2\sin x \cos x+1=0$.\\
		Đặt $t=\sin x+\cos x=\sqrt{2} \sin \left(x+\dfrac{\pi}{4}\right) \Rightarrow \sin x \cos x=\dfrac{t^2-1}{2}$, điều kiện $-\sqrt{2} \leq t \leq \sqrt{2}$.\\
		Phương trình trở thành: $2 t+\left(t^2-1\right)+1=0\Leftrightarrow t^2+2 t=0\Leftrightarrow\hoac{&t=0\quad (\text{nhận}) \\& t=-2\quad(\text {loại }).}$\\
		Khi đó $\sin x+\cos x=0\Leftrightarrow \sin \left(x+\dfrac{\pi}{4}\right)=0\Leftrightarrow x=-\dfrac{\pi}{4}+k \pi, \quad k \in \mathbb{Z}$.
		\\
		Vậy phương trình có nghiệm $x=-\dfrac{\pi}{4}+k \pi, \quad k \in \mathbb{Z}$.}
\end{bt}

% nguồn: 2026-11-15_11_1D1_TL2_Luong_giac (nhập 15/11/2026)
\begin{bt}%[1D1V6-4]%[ma:161ce420]
% nguon: Archive/SP_KhoLatex/ID6_old/1K1-4-10.tex #3615
%[DCHT-11-KNTT]%[Vĩ Lê Văn]
	Giải phương trình	$
	\sin x \cos x=6(\sin x-\cos x-1)$.
		\dapso{$\hoac{&x=\dfrac{\pi}{2}+k \pi \\& x=\pi+k 2\pi},\quad k \in \mathbb{Z}$.}
		\loigiai{
		Đặt $t=\sin x-\cos x=\sqrt{2} \sin \left(x-\dfrac{\pi}{4}\right) \Rightarrow \sin x \cos x=\dfrac{1-t^2}{2}$, điều kiện $-\sqrt{2} \leq t \leq \sqrt{2}$.\\
		Phương trình trở thành: $1-t^2=12(t-1) \Leftrightarrow t^2+12 t-13=0\Leftrightarrow\hoac{&t=1 \quad (\text{nhận}) \\& t=-13\quad(\text {loại}).}$\\
		Vớt $t=1$, ta có $\sin x-\cos x=1\Leftrightarrow \sin \left(x-\dfrac{\pi}{4}\right)=\dfrac{1}{\sqrt{2}}\Leftrightarrow \hoac{&x=\dfrac{\pi}{2}+k \pi \\& x=\pi+k 2\pi}, \quad k \in \mathbb{Z}$.\\
		Vậy phương trình có nghiệm: $\hoac{&x=\dfrac{\pi}{2}+k \pi \\& x=\pi+k 2\pi},\quad k \in \mathbb{Z}$.}
	
\end{bt}

% nguồn: 2026-11-15_11_1D1_TL2_Luong_giac (nhập 15/11/2026)
\begin{bt}%[1D1V6-4]%[ma:8c31c47a]
% nguon: Archive/SP_KhoLatex/ID6_old/1K1-4-10.tex #3617
%[DCHT-11-KNTT]%[Vĩ Lê Văn]
	Giải phương trình $\sin 2 x+\sqrt{2} \sin \left(x-\dfrac{\pi}{4}\right)=1$.
		\dapso{$\hoac{& x=\dfrac{\pi}{4}+k \pi\\&x=\dfrac{\pi}{2}+k 2\pi \\& x=\pi+k 2\pi}, \quad k \in \mathbb{Z}$.}
	\loigiai{
		Ta có
		$\sin 2 x+\sqrt{2} \sin \left(x-\dfrac{\pi}{4}\right)=1\Leftrightarrow 2\sin x \cos x+\sin x-\cos x=1$.\\
		Đặt $t=\sin x-\cos x=\sqrt{2} \sin \left(x-\dfrac{\pi}{4}\right) \Rightarrow \sin x \cos x=\dfrac{1-t^2}{2}$, điều kiện $-\sqrt{2} \leq t \leq \sqrt{2}$.\\
		Phương trình trở thành: $1-t^2+t=1\Leftrightarrow t^2-t=0\Leftrightarrow\hoac{&t=0\quad (\text{nhận})\\& t=1\quad (\text{nhận})}$\\
		$\Leftrightarrow \hoac{&\sqrt{2}sin \left(x-\dfrac{\pi}{4}\right)=0\\&\sqrt{2}sin \left(x-\dfrac{\pi}{4}\right)=1}\Leftrightarrow \hoac{&sin \left(x-\dfrac{\pi}{4}\right)=0\\&\sin \left(x-\dfrac{\pi}{4}\right)=\dfrac{1}{\sqrt{2}}}\Leftrightarrow \hoac{& x=\dfrac{\pi}{4}+k \pi\\&x=\dfrac{\pi}{2}+k 2\pi \\& x=\pi+k 2\pi}, \quad k \in \mathbb{Z}$.\\
		Vậy phương trình có nghiệm: $\hoac{& x=\dfrac{\pi}{4}+k \pi\\&x=\dfrac{\pi}{2}+k 2\pi \\& x=\pi+k 2\pi}, \quad k \in \mathbb{Z}$.}
\end{bt}

% nguồn: 2026-11-15_11_1D1_TL2_Luong_giac (nhập 15/11/2026)
\begin{bt}%[1D1V6-4]%[ma:019c4502]
% nguon: Archive/SP_KhoLatex/ID6_old/1K1-4-10.tex #3618
%[DCHT-11-KNTT]%[Vĩ Lê Văn]
	Giải phương trình  $\sin x-\cos x+2\sin 2x+1=0$.
		\dapso{$\hoac{&x=k2\pi \\&x=\dfrac{3\pi}{2}+k2\pi}, \quad k \in \mathbb{Z}$.}
	\loigiai{
		Đặt $t=\sin x-\cos x=\sqrt{2}\sin\left(x- \dfrac{\pi}{4}\right)$ thì $t \in \left[-\sqrt{2};\sqrt{2}\right]$.\\
		$\Rightarrow \left(\sin x-\cos x\right)^2=t^2 \Leftrightarrow \sin^2 x -2\sin x\cos x+ \cos^2 x =t^2 \Leftrightarrow 1-\sin 2x=t^2 \Leftrightarrow \sin 2x=1-t^2$.\\
		Phương trình trở thành: $t+2(1-t^2)+1=0 \Leftrightarrow -2t^2+t+3=0 \Leftrightarrow \hoac{&t=-1\\&t=\dfrac{3}{2} \quad (\text{loại}).}$\\
		Với $t=-1 \Leftrightarrow \sqrt{2}\sin\left(x- \dfrac{\pi}{4}\right)=-1 \Leftrightarrow \sin \left(x-\dfrac{\pi}{4}\right) =\dfrac{-1}{\sqrt{2}} \Leftrightarrow \hoac{&x=k2\pi \\&x=\dfrac{3\pi}{2}+k2\pi}, \quad k \in \mathbb{Z}.$\\
		Vậy nghiệm của phương trình là $\hoac{&x=k2\pi \\&x=\dfrac{3\pi}{2}+k2\pi}, \quad k \in \mathbb{Z}$.}

\end{bt}

% nguồn: 2026-11-15_11_1D1_TL2_Luong_giac (nhập 15/11/2026)
\begin{bt}%[1D1V6-4]%[ma:fcc2c36e]
% nguon: Archive/SP_KhoLatex/ID6_old/1K1-4-10.tex #3619
%[DCHT-11-KNTT]%[Vĩ Lê Văn]
	Giải phương trình  $\sin x+\cos x+\sin x\cos x=1$.
		\dapso{$\hoac{&x=k2\pi \\&x=\dfrac{\pi}{2}+k2\pi}, \quad k \in \mathbb{Z}$.}
	\loigiai{
		Đặt $t=\sin x+\cos x=\sqrt{2}\sin\left(x+ \dfrac{\pi}{4}\right)$ thì $t \in \left[-\sqrt{2};\sqrt{2}\right]$.\\
		$\Rightarrow \left(\sin x+\cos x\right)^2=t^2 \Leftrightarrow 1+2\sin x\cos x=t^2 \Leftrightarrow \sin x\cos x=\dfrac{t^2-1}{2}$.\\
		Phương trình trở thành: $t+\dfrac{t^2-1}{2}=1 \Leftrightarrow t^2+2t-3=0 \Leftrightarrow \hoac{&t=1\\&t=-3 \quad (\text{loại}).}$\\
		Với $t=1 \Leftrightarrow \sqrt{2}\sin\left(x+ \dfrac{\pi}{4}\right)=1 \Leftrightarrow \sin \left(x+\dfrac{\pi}{4}\right) =\dfrac{1}{\sqrt{2}} \Leftrightarrow \hoac{&x=k2\pi \\&x=\dfrac{\pi}{2}+k2\pi}, \quad k \in \mathbb{Z}.$\\
		Vậy nghiệm của phương trình là $\hoac{&x=k2\pi \\&x=\dfrac{\pi}{2}+k2\pi}, \quad k \in \mathbb{Z}$.}

\end{bt}

% nguồn: 2026-11-15_11_1D1_TL2_Luong_giac (nhập 15/11/2026)
\begin{bt}%[1D1V6-4]%[ma:0d3edba7]
% nguon: Archive/SP_KhoLatex/ID6_old/1K1-4-10.tex #3620
%[DCHT-11-KNTT]%[Vĩ Lê Văn]
	Giải phương trình  $\sin 2x-2\sqrt{2}\left( \sin x+\cos x\right)=5$.
		\dapso{$x=\dfrac{-3\pi}{4}+k2\pi, \quad k \in \mathbb{Z}$.}
	\loigiai{
		Đặt $t=\sin x+\cos x=\sqrt{2}\sin\left(x+ \dfrac{\pi}{4}\right)$ thì $t \in \left[-\sqrt{2};\sqrt{2}\right]$.\\
		$\Rightarrow \left(\sin x+\cos x\right)^2=t^2 \Leftrightarrow 1+2\sin x\cos x=t^2 \Leftrightarrow \sin 2x=t^2-1$.\\
		Phương trình trở thành: $t^2-1-2\sqrt{2}t=5 \Leftrightarrow t^2-2\sqrt{2}t-6=0 \Leftrightarrow \hoac{&t=-\sqrt{2}\quad (\text{nhận})\\&t=3\sqrt{2} \quad (\text{loại}).}$\\
		Với $t=-\sqrt{2} \Leftrightarrow \sqrt{2}\sin\left(x+ \dfrac{\pi}{4}\right)=-\sqrt{2} \Leftrightarrow \sin \left(x+\dfrac{\pi}{4}\right) =-1 \Leftrightarrow x=\dfrac{-3\pi}{4}+k2\pi, \quad k \in \mathbb{Z}$.\\
		Vậy nghiệm của phương trình là $x=\dfrac{-3\pi}{4}+k2\pi, \quad k \in \mathbb{Z}$.}

\end{bt}

% nguồn: 2026-11-15_11_1D1_TL2_Luong_giac (nhập 15/11/2026)
\begin{bt}%[1D1V6-4]%[ma:1bec6cfa]
% nguon: Archive/SP_KhoLatex/ID6_old/1K1-4-10.tex #3621
%[DCHT-11-KNTT]%[Vĩ Lê Văn]
	Giải phương trình  $5\sin 2x+12=12\left( \sin x-\cos x\right) $.
		\dapso{$\hoac{&x=\pi+k2\pi \\&x=\dfrac{\pi}{2}+k2\pi}, \quad k \in \mathbb{Z}$.}
	\loigiai{
		Đặt $t=\sin x-\cos x=\sqrt{2}\sin\left(x- \dfrac{\pi}{4}\right)$ thì $t \in \left[-\sqrt{2};\sqrt{2}\right]$.\\
		$\Rightarrow \left(\sin x-\cos x\right)^2=t^2 \Leftrightarrow \sin^2 x -2\sin x\cos x \cos^2 x =t^2 \Leftrightarrow 1-\sin 2x=t^2 \Leftrightarrow \sin 2x=1-t^2$.\\
		Phương trình trở thành: $5(1-t^2)+12=12t \Leftrightarrow -5t^2-12t+17=0 \Leftrightarrow \hoac{&t=1\\&t=-\dfrac{17}{5} \quad (\text{loại}).}$\\
		Với $t=1 \Leftrightarrow \sin x -\cos x=1 \Leftrightarrow \sin \left(x-\dfrac{\pi}{4}\right) =\dfrac{1}{\sqrt{2}} \Leftrightarrow \hoac{&x=\pi+k2\pi \\&x=\dfrac{\pi}{2}+k2\pi}, \quad k \in \mathbb{Z}.$\\
		Vậy nghiệm của phương trình là $\hoac{&x=\pi+k2\pi \\&x=\dfrac{\pi}{2}+k2\pi}, \quad k \in \mathbb{Z}$.}
	
\end{bt}

% nguồn: 2026-11-15_11_1D1_TL2_Luong_giac (nhập 15/11/2026)
\begin{bt}%[1D1V6-4]%[ma:13183c05]
% nguon: Archive/SP_KhoLatex/ID6_old/1K1-4-10.tex #3622
%[DCHT-11-KNTT]%[Vĩ Lê Văn]
	Giải phương trình  $\sin x\cos x=6\left( \sin x-\cos x+1\right) $.
		\dapso{$\hoac{&x=k2\pi \\&x=\dfrac{3\pi}{2}+k2\pi}, \quad k \in \mathbb{Z}$.}
	\loigiai{
		Đặt $t=\sin x-\cos x=\sqrt{2}\sin\left(x- \dfrac{\pi}{4}\right)$ thì $t \in \left[-\sqrt{2};\sqrt{2}\right]$.\\
		$\Rightarrow \left(\sin x-\cos x\right)^2=t^2 \Leftrightarrow 1-2\sin x\cos x=t^2 \Leftrightarrow \sin x\cos x=\dfrac{1-t^2}{2}$.\\
		Phương trình trở thành: $\dfrac{1-t^2}{2}=6(t+1) \Leftrightarrow t^2+12t+11=0 \Leftrightarrow \hoac{&t=1\\&t=-11 \quad (\text{loại}).}$\\
		Với $t=-1 \Leftrightarrow \sqrt{2}\sin\left(x- \dfrac{\pi}{4}\right)=-1 \Leftrightarrow \sin \left(x-\dfrac{\pi}{4}\right) =-\dfrac{1}{\sqrt{2}} \Leftrightarrow \hoac{&x=k2\pi \\&x=\dfrac{3\pi}{2}+k2\pi}, \quad k \in \mathbb{Z}$.\\
		Vậy nghiệm của phương trình là $\hoac{&x=k2\pi \\&x=\dfrac{3\pi}{2}+k2\pi}, \quad k \in \mathbb{Z}$.}

\end{bt}

% nguồn: 2026-11-15_11_1D1_TL2_Luong_giac (nhập 15/11/2026)
\begin{bt}%[1D1V6-4]%[ma:3508de69]
% nguon: Archive/SP_KhoLatex/ID6_old/1K1-4-10.tex #3623
%[DCHT-11-KNTT]%[Vĩ Lê Văn]
	Giải phương trình  $\sin x\cos x+\left| \sin x+\cos x\right| =1$.
		\dapso{$x=k\dfrac{\pi}{2}, \quad k \in \mathbb{Z}$.}
	\loigiai{
		Đặt $t=\left| \sin x+\cos x\right|=\left| \sqrt{2}\sin\left(x+ \dfrac{\pi}{4}\right)\right|$ thì $t \in \left[0;\sqrt{2}\right]$.\\
		$\Rightarrow \left|\sin x+\cos x\right|^2=t^2 \Leftrightarrow 1+2\sin x\cos x=t^2 \Leftrightarrow \sin x\cos x=\dfrac{t^2-1}{2}$.\\
		Phương trình trở thành: $\dfrac{t^2-1}{2}+t=1 \Leftrightarrow t^2+2t-3=0 \Leftrightarrow \hoac{&t=1\\&t=-3 \quad (\text{loại}).}$\\
		Với $t=1 \Leftrightarrow \left| \sqrt{2}\sin\left(x+ \dfrac{\pi}{4}\right)\right|=1 \Leftrightarrow \left| \sin \left(x+\dfrac{\pi}{4}\right)\right| =\dfrac{1}{\sqrt{2}} \Leftrightarrow \hoac{&\sin\left(x+\dfrac{\pi}{4}\right)=\dfrac{\sqrt{2}}{2}\\&\sin\left(x+\dfrac{\pi}{4}\right)=-\dfrac{\sqrt{2}}{2}} $\\
		$\Leftrightarrow \hoac{&x=k2\pi \\&x=\dfrac{\pi}{2}+k2\pi\\&x=-\dfrac{\pi}{2}+k2\pi\\&x=\pi+k2\pi} \Leftrightarrow x=k\dfrac{\pi}{2}, \quad k \in \mathbb{Z}.$\\
		Vậy nghiệm của phương trình là $x=k\dfrac{\pi}{2}, \quad k \in \mathbb{Z}$.}

\end{bt}

% nguồn: 2026-11-15_11_1D1_TL2_Luong_giac (nhập 15/11/2026)
\begin{bt}%[1D1V6-4]%[ma:705fe55e]
% nguon: Archive/SP_KhoLatex/ID6_old/1K1-4-10.tex #3624
%[DCHT-11-KNTT]%[Vĩ Lê Văn]
	Giải phương trình  $\left| \cos x-\sin x\right|+3\sin 2x =1$.
		\dapso{$x=k\dfrac{\pi}{2}, \quad k \in \mathbb{Z}$.}
	\loigiai{
		Đặt $t=\left| \cos x-\sin x\right|=\left|\sqrt{2}\sin\left(x- \dfrac{\pi}{4}\right)\right|$ thì $t \in \left[0;\sqrt{2}\right]$.\\
		$\Rightarrow \left| \cos x-\sin x\right|^2=t^2 \Leftrightarrow 1-2\sin x\cos x=t^2 \Leftrightarrow \sin 2x=1-t^2$.\\
		Phương trình trở thành: $t+3(1-t^2)=1 \Leftrightarrow -3t^2+t+2=0 \Leftrightarrow \hoac{&t=1\\&t=-\dfrac{2}{3} \quad (\text{loại}).}$\\
		Với $t=1 \Leftrightarrow \left| \sqrt{2}\sin\left(x- \dfrac{\pi}{4}\right)\right|=1 \Leftrightarrow \left| \sin \left(x-\dfrac{\pi}{4}\right)\right| =\dfrac{1}{\sqrt{2}} \Leftrightarrow \hoac{&\sin\left(x-\dfrac{\pi}{4}\right)=\dfrac{\sqrt{2}}{2}\\&\sin\left(x-\dfrac{\pi}{4}\right)=-\dfrac{\sqrt{2}}{2}}$\\
		$\Leftrightarrow \hoac{&x=k2\pi \\&x=\dfrac{\pi}{2}+k2\pi\\&x=-\dfrac{\pi}{2}+k2\pi\\&x=-\pi+k2\pi}, \quad k \in \mathbb{Z} \Leftrightarrow x=k\dfrac{\pi}{2}, \quad k \in \mathbb{Z}$.\\
		Vậy nghiệm của phương trình là $x=k\dfrac{\pi}{2}, \quad k \in \mathbb{Z}$.}
\end{bt}

% nguồn: 2026-11-15_11_1D1_TL2_Luong_giac (nhập 15/11/2026)
\begin{bt}%[1D1V6-4]%[ma:68510eec]
% nguon: Archive/SP_KhoLatex/ID6_old/1K1-4-10.tex #3625
%[DCHT-11-KNTT]%[Vĩ Lê Văn]
	Giải phương trình: $\tan^2x+\cot^2x-(\tan x-\cot x)-2=0$.
	\dapso{$\hoac{&x=\dfrac{\pi}{4}+k\dfrac{\pi}{2}\\&x=\arctan\dfrac{1+\sqrt{5}}{2}+k\pi\\&x=\arctan\dfrac{1-\sqrt{5}}{2}+k\pi}, \quad k \in \mathbb{Z}$.}
	\loigiai{
		Điều kiện: $x \neq \dfrac{k\pi}{2} \; (k \in \mathbb{Z})$. Đặt $t=\tan x-\cot x$.\\
		$\Rightarrow \left( \tan x-\cot x\right) ^2=t^2 \Leftrightarrow \tan^2x-2\tan x\cot x+\cot^2x=t^2 \Leftrightarrow \tan^2x+\cot^2x=t^2+2$.\\
		Phương trình trở thành: $t^2+2-t-2=0 \Leftrightarrow t^2-t=0 \Leftrightarrow \hoac{&t=0\\&t=1.}$\\
		Với $t=0 \Leftrightarrow \tan x-\cot x =0 \Leftrightarrow \tan x =\tan \left(\dfrac{\pi}{2}-x\right) \Leftrightarrow x=\dfrac{\pi}{2}-x+k\pi \Leftrightarrow x=\dfrac{\pi}{4}+k\dfrac{\pi}{2}, \quad k \in \mathbb{Z}$.\\
		Với  $t=1 \Leftrightarrow \tan x-\cot x =1 \Leftrightarrow \tan x -\dfrac{1}{\tan x}=1 \Leftrightarrow \tan^2x-\tan x -1 =0$\\
		$\Leftrightarrow \hoac{&\tan x=\dfrac{1+\sqrt{5}}{2}\\&\tan x=\dfrac{1-\sqrt{5}}{2}} \Leftrightarrow \hoac{&x=\arctan\dfrac{1+\sqrt{5}}{2}+k\pi \\&x=\arctan\dfrac{1-\sqrt{5}}{2}+k\pi}, \quad k \in \mathbb{Z}$.\\
		Vậy nghiệm của phương trình là $\hoac{&x=\dfrac{\pi}{4}+k\dfrac{\pi}{2}\\&x=\arctan\dfrac{1+\sqrt{5}}{2}+k\pi\\&x=\arctan\dfrac{1-\sqrt{5}}{2}+k\pi}, \quad k \in \mathbb{Z}$.}
	
\end{bt}

% nguồn: 2026-11-15_11_1D1_TL2_Luong_giac (nhập 15/11/2026)
\begin{bt}%[1D1V6-4]%[ma:43aa0970]
% nguon: Archive/SP_KhoLatex/ID6_old/1K1-4-10.tex #3626
%[DCHT-11-KNTT]%[Vĩ Lê Văn]
	Giải phương trình: $3\tan^2x+4\tan x+4\cot x+3\cot^2x+2=0$.
	\dapso{$x=-\dfrac{\pi}{4}+k\pi, \quad k \in \mathbb{Z}$.}
	\loigiai{
		Điều kiện: $x \neq \dfrac{k\pi}{2} \; (k \in \mathbb{Z})$. Đặt $t=\tan x+\cot x$.\\
		$\Rightarrow \left( \tan x+\cot x\right) ^2=t^2 \Leftrightarrow \tan^2x+2\tan x\cot x+\cot^2x=t^2 \Leftrightarrow \tan^2x+\cot^2x=t^2-2$\\
		Phương trình trở thành: $3(t^2-2)+4t+2=0 \Leftrightarrow 3t^2+4t-4=0 \Leftrightarrow \hoac{&t=-2\\&t=\dfrac{2}{3}.}$\\
		Với  $t=-2 \Leftrightarrow \tan x+\cot x =-2 \Leftrightarrow \tan x +\dfrac{1}{\tan x}=-2 \Leftrightarrow \tan^2x+2\tan x +1 =0$\\
		$\Leftrightarrow \tan x=-1 \Leftrightarrow x=-\dfrac{\pi}{4}+k\pi, \quad k \in \mathbb{Z}$.\\
		Với $t=\dfrac{2}{3} \Leftrightarrow \tan x+\cot x =\dfrac{2}{3} \Leftrightarrow \tan x +\dfrac{1}{\tan x}=\dfrac{2}{3} \Leftrightarrow \tan^2x-\dfrac{2}{3}\tan x +1 =0$.\\
		Phương trình này vô nghiệm.\\
		Vậy nghiệm của phương trình là $x=-\dfrac{\pi}{4}+k\pi, \quad k \in \mathbb{Z}$.}

\end{bt}
